\section{Introduction}
\label{sec:intro}

Dense matrix multiplication (GEMM) is the standard workload for testing a compute accelerator end to end: it needs
high arithmetic throughput, it stresses the memory system unless it is tiled well, and its result is easy to check.
This project puts a GEMM accelerator on the \textbf{DU-PCB-001-003}, a PCIe card with a Kintex-7 480T FPGA and two
2\,GB DDR3 channels, using open tools around the vendor flow:
\begin{itemize}
  \item \textbf{LiteX}~\cite{litex} with LitePCIe~\cite{litepcie} and LiteDRAM~\cite{litedram} for the SoC: PCIe
        endpoint and DMA, both DDR3 controllers, the control registers and the glue to the kernel;
  \item \textbf{Vitis HLS}~\cite{ug1399} for the kernel, a copy of the open-source \code{gemm\_hls}~\cite{gemmhls,gemmhlsrepo}
        adapted to run without a vendor shell;
  \item \textbf{Vivado} for implementation, driven from LiteX's build script with our Tcl hooks;
  \item a \textbf{C++ host program} on the LitePCIe Linux driver, and a \textbf{Python flow} that builds, loads,
        runs and explores configurations from JSON files.
\end{itemize}

\paragraph{What the system does.} The host sends \(A\) (as \(A^\mathsf{T}\)) and \(B\) over PCIe DMA into the DDR3
channels, starts the kernel through control registers, waits for it, and reads \(C\) back together with performance
counters: kernel cycles, per-matrix AXI traffic, the cycles each read channel waited for data, and read latencies.
Every result is verified against a host reference and written as one JSON row with the build id, the configuration
hash and the git state, so any number in this report can be traced to the exact bitstream and source that produced it.

\paragraph{Contributions.}
\begin{enumerate}
  \item A CPU-less LiteX SoC with a PCIe-DMA-to-DDR3 bridge for both channels and an HLS kernel wired to LiteDRAM
        through AXI, with per-bundle performance counters (\cref{sec:system,sec:pcie,sec:ddr,sec:kernel}).
  \item A C++ host runtime on the LitePCIe driver, with zero-copy DMA rings and a fix that doubles card-to-host
        throughput (\cref{sec:pcie,sec:host}).
  \item A JSON-driven, cached build flow and a DSE driver that runs whole campaigns unattended and keeps every
        result traceable (\cref{sec:flow,sec:dse}).
  \item A measured exploration over processing elements, data layout, tile sizes, memory maps, matrix sizes, clocks
        and Vivado strategies, ending at \gflops{110.9} (\cref{sec:eval}).
  \item Integer and fixed-point formats through the same kernel source, with exact verification on the card, up to
        \SI{1370.9}{GOP/s} for int4; and two silent-error causes found on the way, DDR3 read margin under load and
        unchecked tile rules in the kernel (\cref{sec:formats}).
  \item DenseShift power-of-two weights from training to the card: PyTorch layers, a mixed-type kernel with a shift
        MAC (3.4\x\ the fastest float32\,\(\times\)\,float32 throughput with int8 activations and 40 DSPs, 1.7\x\
        with float activations) and an MNIST network whose layers run on the accelerator 3.1\x\ faster than in
        float32, with results identical to the CPU (\cref{sec:ds}).
\end{enumerate}

\paragraph{Results at a glance.} \Cref{tab:headline} lists the configurations that mark the path to the best result.
The full exploration, with the reasons for each step, is in \cref{sec:dse,sec:eval}. The best build per integer and
fixed-point format is in \cref{tab:formats} (\cref{sec:formats}).

\begin{table}[ht]
\centering\small
\caption{Selected configurations (float32, kernel time, all results verified). Peak = \(2 P_N P_M f\).}
\label{tab:headline}
\begin{tabular}{@{}llrrrl@{}}
\toprule
Config & Clock & Peak & Best & \% & Resources (LUT / DSP / BRAM tiles) \\
\midrule
PE 8\x8, tile 128 (first build) & 200\,MHz & 25.6 & 23.4 & 91.5 & 64.6\,k / 383 / 238.5 \\
PE 8\x8, tile 128 & 250\,MHz & 32.0 & 29.0 & 90.7 & 67.1\,k / 383 / 231 \\
PE 16\x8, tile 256 & 200\,MHz & 51.2 & 47.2 & 92.2 & 94.6\,k / 703 / 303 \\
PE 32\x8, tile 256 & 200\,MHz & 102.4 & 86.5 & 84.5 & 154.2\,k / 1343 / 414.5 \\
PE 32\x8, A\(^\mathsf{T}\), tile 1024\x512 & 200\,MHz & 102.4 & 98.6 & 96.3 & 154.3\,k / 1320 / 654.5 \\
\textbf{PE 32\x8, A\(^\mathsf{T}\), tile 1024\x512, aggressive} & \textbf{225\,MHz} & 115.2 & \textbf{110.9} & 96.2 & 152.0\,k / 1320 / 654.5 \\
\bottomrule
\end{tabular}
\end{table}
